\documentclass[a4paper,landscape,8pt]{article}
\usepackage[utf8]{inputenc}
\usepackage[T1]{fontenc}
\usepackage[spanish]{babel}
\usepackage{geometry}
\usepackage{array}
\usepackage{colortbl}
\usepackage{xcolor}
\usepackage{booktabs}
\usepackage{longtable}
\usepackage{fancyhdr}
\usepackage{pdflscape}

\geometry{left=1cm, right=1cm, top=1.8cm, bottom=1.5cm}

% Colores
\definecolor{violeta}{HTML}{7B2D8E}
\definecolor{naranja}{HTML}{E67E22}
\definecolor{verde}{HTML}{27AE60}
\definecolor{marron}{HTML}{8B4513}
\definecolor{azulclaro}{HTML}{2980B9}
\definecolor{amarillo}{HTML}{D4AC0D}
\definecolor{azulg}{HTML}{3498DB}
\definecolor{rojog}{HTML}{E74C3C}
\definecolor{gris}{HTML}{F2F2F2}

\pagestyle{fancy}
\fancyhf{}
\fancyhead[L]{\small\textbf{Nuevo Compañeros 1 — U03: La Familia}}
\fancyhead[R]{\small\textbf{Tarjetas de vocabulario — Caja 1}}
\fancyfoot[C]{\thepage}

\setlength{\parindent}{0pt}
\renewcommand{\arraystretch}{1.25}

\newcommand{\M}{\textcolor{azulg}{M}}
\newcommand{\F}{\textcolor{rojog}{F}}
\newcommand{\MF}{M/F}

\begin{document}

\begin{center}
{\Large\textbf{TARJETAS DE VOCABULARIO — UNIDAD 3: LA FAMILIA}}\\[0.2cm]
{\normalsize Sección Vocabulario (p.34--35) — Caja 1: Estación de servicio}
\end{center}

\vspace{0.3cm}

%% ============================================================
%% CAJA 1 — FAMILIA (violeta) — 18 tarjetas
%% ============================================================

\noindent\textcolor{violeta}{\rule{\linewidth}{2pt}}
\section*{\textcolor{violeta}{FAMILIA} {\small — 18 tarjetas — campo violeta — Frec: 3=alta, 2=media, 1=baja}}

{\footnotesize
\begin{longtable}{|>{\bfseries}p{1.4cm}|p{0.6cm}|p{0.7cm}|p{1.5cm}|p{3.2cm}|p{3.5cm}|p{0.5cm}|p{2.5cm}|p{1cm}|p{1cm}|p{0.8cm}|p{1.1cm}|p{1cm}|p{1.2cm}|}
\hline
\rowcolor{violeta!15}
Palabra & Gén. & Color & Sílaba & Regla & Ejemplo & Fr. & Irregularidad & IT & FR & PT-BR & EN & CS & PL \\
\hline
\endfirsthead
\hline
\rowcolor{violeta!15}
Palabra & Gén. & Color & Sílaba & Regla & Ejemplo & Fr. & Irregularidad & IT & FR & PT-BR & EN & CS & PL \\
\hline
\endhead

padre & \M & azul & PA-dre & Irregular: raíces distintas (padre/madre). & El padre de Javier es camionero. & 3 & Raíces distintas (padre/madre). & padre & père & pai & father & otec & ojciec \\
\hline
madre & \F & rojo & MA-dre & Irregular: raíces distintas (padre/madre). & La madre de Javier se llama Catalina y es enfermera. & 3 & Raíces distintas (madre/madra). & madre & mère & mãe & mother & matka & matka \\
\hline
\rowcolor{gris}
\multicolumn{14}{|l|}{\textit{Plural mixto: padre + madre = \textbf{padres} (parents)}} \\
\hline

hermano & \M & azul & her-MA-no & -o/-a: hermano/hermana & Carlos es el hermano mayor de Lucía. & 3 & & fratello & frère & irmão & brother & bratr & brat \\
\hline
hermana & \F & rojo & her-MA-na & -o/-a: hermano/hermana & Javier tiene una hermana pequeña, Alejandra. & 3 & & sorella & s\oe{}ur & irmã & sister & sestra & siostra \\
\hline
\rowcolor{gris}
\multicolumn{14}{|l|}{\textit{Plural mixto: hermano + hermana = \textbf{hermanos} (siblings)}} \\
\hline

hijo & \M & azul & HI-jo & -o/-a: hijo/hija & David es el hijo de Alicia y Luis. & 2 & & figlio & fils & filho & son & syn & syn \\
\hline
hija & \F & rojo & HI-ja & -o/-a: hijo/hija & María es la hija de Carlos y Carmen. & 2 & & figlia & fille & filha & daughter & dcera & córka \\
\hline
\rowcolor{gris}
\multicolumn{14}{|l|}{\textit{Plural mixto: hijo + hija = \textbf{hijos} (children)}} \\
\hline

abuelo & \M & azul & a-BUE-lo & -o/-a: abuelo/abuela & El abuelo de David se llama Carlos. & 2 & & nonno & grand-père & avô & grandfather & dědeček & dziadek \\
\hline
abuela & \F & rojo & a-BUE-la & -o/-a: abuelo/abuela & La abuela de David se llama Carmen. & 2 & & nonna & grand-mère & avó & grandmother & babička & babcia \\
\hline
\rowcolor{gris}
\multicolumn{14}{|l|}{\textit{Plural mixto: abuelo + abuela = \textbf{abuelos} (grandparents)}} \\
\hline

tío & \M & azul & TÍ-o & -o/-a: tío/tía & Nacho es el tío de David. & 2 & & zio & oncle & tio & uncle & strýc & wujek \\
\hline
tía & \F & rojo & TÍ-a & -o/-a: tío/tía & María es la tía de David. & 2 & & zia & tante & tia & aunt & teta & ciotka \\
\hline
\rowcolor{gris}
\multicolumn{14}{|l|}{\textit{Plural mixto: tío + tía = \textbf{tíos} (uncle and aunt)}} \\
\hline

primo & \M & azul & PRI-mo & -o/-a: primo/prima & Álvaro es el primo de David. & 2 & & cugino & cousin & primo & cousin & bratranec & kuzyn \\
\hline
prima & \F & rojo & PRI-ma & -o/-a: primo/prima & Paloma es la prima de David. & 2 & & cugina & cousine & prima & cousin & sestřenice & kuzynka \\
\hline

marido & \M & azul & ma-RI-do & Irregular: par marido/mujer. & Luis es el marido de Alicia. & 2 & Par irregular: marido/mujer. & marito & mari & marido & husband & manžel & mąż \\
\hline
mujer & \F & rojo & mu-JER & Irregular: par marido/mujer. & Alicia es la mujer de Luis. & 3 & Doble significado: woman / wife. & moglie & femme & esposa & wife & manželka & żona \\
\hline

nieto & \M & azul & NIE-to & -o/-a: nieto/nieta & David es el nieto de Carlos y Carmen. & 1 & & nipote & petit-fils & neto & grandson & vnuk & wnuk \\
\hline
nieta & \F & rojo & NIE-ta & -o/-a: nieto/nieta & Pilar es la nieta de Carlos y Carmen. & 1 & & nipote & petite-fille & neta & grand-daughter & vnučka & wnuczka \\
\hline
\rowcolor{gris}
\multicolumn{14}{|l|}{\textit{Plural mixto: nieto + nieta = \textbf{nietos} (grandchildren)}} \\
\hline

sobrino & \M & azul & so-BRI-no & -o/-a: sobrino/sobrina & David es el sobrino de Roberto y Carmen. & 1 & & nipote & neveu & sobrinho & nephew & synovec & siostrzeniec \\
\hline
sobrina & \F & rojo & so-BRI-na & -o/-a: sobrino/sobrina & Paloma es la sobrina de Alicia y Luis. & 1 & & nipote & nièce & sobrinha & niece & neteř & siostrzenica \\
\hline

\end{longtable}
}

\newpage

%% ============================================================
%% CAJA 1 — PROFESIONES (naranja) — 3 tarjetas (textos p.35)
%% ============================================================

\noindent\textcolor{naranja}{\rule{\linewidth}{2pt}}
\section*{\textcolor{naranja}{PROFESIONES} {\small — 3 tarjetas — campo naranja — Textos p.35 (Javier y Lucía)}}

{\footnotesize
\begin{longtable}{|>{\bfseries}p{1.6cm}|p{0.6cm}|p{0.7cm}|p{1.8cm}|p{3.5cm}|p{3.5cm}|p{0.5cm}|p{2.8cm}|p{1.2cm}|p{1.3cm}|p{1.3cm}|p{1.2cm}|p{1.5cm}|p{1.5cm}|}
\hline
\rowcolor{naranja!15}
Palabra & Gén. & Color & Sílaba & Regla & Ejemplo & Fr. & Irregularidad & IT & FR & PT-BR & EN & CS & PL \\
\hline

camionero & \M & azul & ca-mio-NE-ro & -o/-a: camionero/camionera & El padre de Javier es camionero. & 1 & & camionista & camionneur & caminhoneiro & truck driver & řidič kamionu & kierowca \\
\hline
enfermera & \F & rojo & en-fer-ME-ra & -o/-a: enfermero/enfermera & La madre de Javier, Catalina, es enfermera. & 1 & & infermiera & infirmière & enfermeira & nurse & zdravotní sestra & pielęgniarka \\
\hline
agricultor & \M & azul & a-gri-cul-TOR & -tor/-tora: agricultor/agricultora. No sigue -o/-a. & El padre de Lucía es agricultor. & 1 & No sigue -o/-a: -tor/-tora (como doctor/doctora). & agricoltore & agriculteur & agricultor & farmer & zemědělec & rolnik \\
\hline

\end{longtable}
}

\vspace{0.5cm}

%% ============================================================
%% CAJA 1 — LUGARES (verde) — 4 tarjetas nuevas
%% ============================================================

\noindent\textcolor{verde}{\rule{\linewidth}{2pt}}
\section*{\textcolor{verde}{LUGARES} {\small — 4 tarjetas — campo verde — Textos p.35}}

{\footnotesize
\begin{longtable}{|>{\bfseries}p{1.6cm}|p{0.6cm}|p{0.7cm}|p{1.8cm}|p{3.5cm}|p{3.8cm}|p{0.5cm}|p{2.5cm}|p{1cm}|p{1.2cm}|p{1cm}|p{1.5cm}|p{1.5cm}|p{1.2cm}|}
\hline
\rowcolor{verde!15}
Palabra & Gén. & Color & Sílaba & Regla & Ejemplo & Fr. & Irregularidad & IT & FR & PT-BR & EN & CS & PL \\
\hline

hotel & \M & azul & ho-TEL & Masculino. El hotel. & Su madre trabaja en el hotel del pueblo. & 1 & & hotel & hôtel & hotel & hotel & hotel & hotel \\
\hline
pueblo & \M & azul & PUE-blo & Masculino: termina en -o. & Lucía es de un pueblo de Cantabria. & 1 & & paese & village & vilarejo & village & vesnice & wieś \\
\hline
instituto & \M & azul & ins-ti-TU-to & Masculino: termina en -o. & Javier estudia en el instituto de su barrio. & 1 & & istituto & lycée & escola & secondary school & střední škola & liceum \\
\hline
barrio & \M & azul & BA-rrio & Masculino: termina en -o. & El instituto está en su barrio. & 1 & & quartiere & quartier & bairro & neighbourhood & čtvrť & dzielnica \\
\hline

\end{longtable}
}

\vspace{0.5cm}

%% ============================================================
%% CAJA 1 — OCIO Y ACTIVIDADES (azul) — 2 tarjetas nuevas
%% ============================================================

\noindent\textcolor{azulclaro}{\rule{\linewidth}{2pt}}
\section*{\textcolor{azulclaro}{OCIO Y ACTIVIDADES} {\small — 2 tarjetas — campo azul — Textos p.35}}

{\footnotesize
\begin{longtable}{|>{\bfseries}p{1.6cm}|p{0.6cm}|p{0.7cm}|p{1.8cm}|p{3.5cm}|p{3.8cm}|p{0.5cm}|p{2.5cm}|p{1.2cm}|p{1cm}|p{1cm}|p{1cm}|p{1cm}|p{1.5cm}|}
\hline
\rowcolor{azulclaro!15}
Palabra & Gén. & Color & Sílaba & Regla & Ejemplo & Fr. & Irregularidad & IT & FR & PT-BR & EN & CS & PL \\
\hline

fútbol & \M & azul & FÚT-bol & Masculino. El fútbol. & Javier juega al fútbol tres días a la semana. & 1 & & calcio & football & futebol & football & fotbal & piłka nożna \\
\hline
piano & \M & azul & pia-NO & Masculino: termina en -o. & Lucía estudia piano. & 1 & & pianoforte & piano & piano & piano & klavír & fortepian \\
\hline

\end{longtable}
}

\vspace{0.5cm}

%% ============================================================
%% CAJA 1 — ESTUDIOS (amarillo) — 1 tarjeta nueva
%% ============================================================

\noindent\textcolor{amarillo}{\rule{\linewidth}{2pt}}
\section*{\textcolor{amarillo}{ESTUDIOS} {\small — 1 tarjeta — campo amarillo — Texto p.35 (Lucía)}}

{\footnotesize
\begin{longtable}{|>{\bfseries}p{1.6cm}|p{0.6cm}|p{0.7cm}|p{1.8cm}|p{3.5cm}|p{3.8cm}|p{0.5cm}|p{2.5cm}|p{1.2cm}|p{1.3cm}|p{1.3cm}|p{1.3cm}|p{1.5cm}|p{1.3cm}|}
\hline
\rowcolor{amarillo!15}
Palabra & Gén. & Color & Sílaba & Regla & Ejemplo & Fr. & Irregularidad & IT & FR & PT-BR & EN & CS & PL \\
\hline

ingeniería & \F & rojo & in-ge-nie-RÍ-a & Femenino: termina en -ía. & Ana, la hermana mayor, estudia ingeniería. & 1 & & ingegneria & ingénierie & engenharia & engineering & inženýrství & inżynieria \\
\hline

\end{longtable}
}

\vspace{0.5cm}

%% ============================================================
%% CAJA 1 — ANIMALES (marrón) — 1 tarjeta nueva
%% ============================================================

\noindent\textcolor{marron}{\rule{\linewidth}{2pt}}
\section*{\textcolor{marron}{ANIMALES} {\small — 1 tarjeta — campo marrón — Texto p.35 (Lucía)}}

{\footnotesize
\begin{longtable}{|>{\bfseries}p{1.6cm}|p{0.6cm}|p{0.7cm}|p{1.8cm}|p{3.5cm}|p{3.8cm}|p{0.5cm}|p{2.5cm}|p{1cm}|p{1cm}|p{1cm}|p{1cm}|p{1cm}|p{1cm}|}
\hline
\rowcolor{marron!15}
Palabra & Gén. & Color & Sílaba & Regla & Ejemplo & Fr. & Irregularidad & IT & FR & PT-BR & EN & CS & PL \\
\hline

vaca & \F & rojo & VA-ca & Femenino: la vaca. Masc: el toro (raíces distintas). & El padre de Lucía es agricultor y también tiene vacas. & 1 & Vaca/toro: raíces distintas. & mucca & vache & vaca & cow & kráva & krowa \\
\hline

\end{longtable}
}

\vspace{1cm}

%% ============================================================
%% RESUMEN
%% ============================================================

\begin{center}
{\large\textbf{RESUMEN — Caja 1: Vocabulario U03 (p.34--35)}}\\[0.4cm]

\begin{tabular}{|l|c|l|}
\hline
\rowcolor{gris}
\textbf{Campo semántico} & \textbf{Tarjetas} & \textbf{Color} \\
\hline
\textcolor{violeta}{Familia} & 18 & violeta \\
\hline
\textcolor{naranja}{Profesiones} & 3 & naranja \\
\hline
\textcolor{verde}{Lugares} & 4 & verde \\
\hline
\textcolor{azulclaro}{Ocio y actividades} & 2 & azul \\
\hline
\textcolor{amarillo}{Estudios} & 1 & amarillo \\
\hline
\textcolor{marron}{Animales} & 1 & marrón \\
\hline
\textbf{TOTAL} & \textbf{29} & \\
\hline
\end{tabular}
\end{center}

\end{document}
